
\section{General Setup}
\label{sec:general_setup}

\noindent\textbf{Training data \& ID/OOD evaluation.} As stated in \S\ref{sec:intro}, we are interested in whether transformers can induce and apply latent rules over knowledge implicitly in a generalizable way. We create a data-generating process consisting of 1) sampling a set of basic \textit{atomic facts}, and 2) using the atomic facts and latent rules to deduce \textit{inferred facts}. To better characterize the level of generalization acquired by the model, we evaluate the model's in-distribution (ID) and out-of-distribution (OOD) performance.
We prepare two separate sets of atomic facts: \ttsmall{atomic$_{\text{ID}}$} and \ttsmall{atomic$_{\text{OOD}}$}. Our training set includes \textit{all} the atomic facts and a uniformly random portion of the inferred facts deduced from \ttsmall{atomic$_{\text{ID}}$}, which we call \ttsmall{train\_inferred$_{\text{ID}}$}. For evaluation, (1) ID generalization aims to evaluate whether the model learns the latent rules correctly, by testing its ability to complete unseen inferred facts also deduced from \ttsmall{atomic$_{\text{ID}}$}, which we denote by \ttsmall{test\_inferred$_{\text{ID}}$}. (2) OOD generalization aims to evaluate the \textit{systematicity}~\citep{lake2018generalization} acquired by the model, namely, the ability to apply rules over knowledge regardless of its distribution. To do this, we test the model on the facts deduced from \ttsmall{atomic$_{\text{OOD}}$}, denoted by \ttsmall{test\_inferred$_{\text{OOD}}$}.

\noindent\textbf{Model \& optimization.} We use a standard decoder-only transformer model as in GPT-2~\citep{radford2019language} with \num{8} layers, \num{768} hidden dimensions and \num{12} attention heads (we explore the impact of different model scales in Appendix~\ref{app:scaling}). Optimization is done by AdamW~\citep{loshchilov2018decoupled} with learning rate $10^{-4}$, batch size \num{512}, weight decay \num{0.1} and \num{2000} warm-up steps. Notably, models are trained for a large number of epochs/steps beyond the point where training performance saturates. More details are in Appendix~\ref{app:compute}.

